%%%%%%%%%%%%%%%%%%%%%%%%%%%%%%%%%%%%%%%%%%%%%%%%%%%%%%%%%%%%%%%%%%%%%%%%%%%
% Lehreinheit 11: Wissenschaftliches Schreiben
% Quelle: WuZ-Buch, Kap. 15 (LE 11), Kap. 12 Praxis-Werkstatt
%         (ch:werkstatt, Schritt 3 "Storytelling"), Kap. 8
%         Begriffsersetzungen NLP -> evidenzbasiert (ch:nlp)
% Staunmoment: Patter-Umbau in Fachsprache  |  Kernbegriff: Präzision & Struktur
%%%%%%%%%%%%%%%%%%%%%%%%%%%%%%%%%%%%%%%%%%%%%%%%%%%%%%%%%%%%%%%%%%%%%%%%%%%
\nsl{Bis hierher hat das Semester Fragen gestellt, Quellen geprüft, Daten
erhoben und gedeutet. In dieser Einheit geht es darum, das alles so
aufzuschreiben, dass andere es prüfen können. Der Einstieg ist ein
Sprechtext, wie ihn Zauberkünstler verwenden: blumig, suggestiv, voller
Behauptungen, die niemand belegt. Die Gruppe übersetzt ihn gemeinsam in
Fachsprache und merkt dabei, dass diese Übersetzung Regeln folgt. Aus den
Regeln entwickelt die Einheit die Grundlagen wissenschaftlichen Schreibens:
den Aufbau eines Textes (IMRaD), den akademischen Stil, die
Argumentationsstruktur, das Zitieren und die Grenze zum Plagiat. Den
Abschluss bilden zwei Werkzeuge für die Seminararbeit: die Einleitung im
Vierschritt und die Überarbeitung als Methode.
\newpage}
\os{\begin{description}
\item[{\rot\bf Heute}]~\\[-\lblskip]
\bi
  \ite ein Seil, das wächst -- und ein Text, der schrumpft
  \ite Aufbau wissenschaftlicher Texte: IMRaD
  \ite Stil, Argument, Begriff: Präzision als Ehrlichkeit
  \ite Zitieren, Paraphrasieren, Plagiat
  \ite Einleitung im Vierschritt, Überarbeiten als Methode
\ei
\end{description}
\newpage}

\begin{vorbereitung}[Patter-Umbau in Fachsprache]
\textbf{Eine Woche vorher}
\begin{checkliste}
  \item Einen Seileffekt auswählen, bei dem ein kurzes Seil sichtbar länger
    wird (handelsübliche Seil-Effekte aus dem Zauberfachhandel). Methode
    und Einstudierung nach der Anleitung des Herstellers bzw.\ der
    Standardliteratur; hier bewusst nicht beschrieben.
  \item Den Effekt mindestens fünfmal mit dem Sprechtext unten proben,
    davon einmal vor einer Testperson. Der Sprechtext ist absichtlich
    übertrieben; er soll ernst, aber nicht ironisch vorgetragen werden,
    sonst nimmt die Gruppe die Analyse vorweg.
  \item Den Ausgangs-Patter (Abschnitt~\ref{sec:le11-patter}) auf eine
    Folie legen und zusätzlich als Text auf dem Übungsblatt
    (Aufgabe~1) bereitstellen.
  \item Die Studierenden per Mail bitten, den aktuellen Stand ihrer
    Seminararbeit (Forschungsfrage, Gliederung, Literaturliste aus LE~3)
    mitzubringen. Ohne diesen Stand ist die Einleitungsübung kaum möglich.
\end{checkliste}

\textbf{Handout und Übungsblatt}
\begin{checkliste}
  \item Handout \enquote{Wissenschaftliches Schreiben} und Übungsblatt
    LE~11 je einmal pro Person kopieren (Druckvorlage
    \texttt{handouts.pdf}, Seiten zu LE~11). Das Übungsblatt ist länger
    als üblich; doppelseitig drucken.
  \item Für den Patter-Umbau je Kleingruppe ein A3-Blatt mit zwei
    Spalten \emph{Sprechtext -- Fachsprache} und einer schmalen dritten
    Spalte \emph{Regel}.
  \item Den schlechten Absatz (Aufgabe~2) zusätzlich je Gruppe einmal in
    großer Schrift (A3) ausdrucken, damit die Gruppe darin markieren kann.
  \item Zwei Farben Marker je Gruppe: Rot für Runde~1 (Inhalt), Blau für
    Runde~2 (Sprache und Form).
\end{checkliste}

\textbf{Am Tag selbst}
\begin{checkliste}
  \item Seil griffbereit, aber verdeckt; die Einheit beginnt mit der
    Vorführung, nicht mit der Begrüßung.
  \item Tafel in zwei Felder teilen: links \enquote{Was sagt der Text?},
    rechts \enquote{Was wissen wir tatsächlich?}. Rechts entstehen die
    Transformationsregeln.
  \item Uhr sichtbar stellen: Die Überarbeitungsübung hat zwei Runden von
    je zehn Minuten, die strikt getrennt bleiben sollen.
\end{checkliste}

\textbf{Risiken und Plan B.} Misslingt der Effekt, hilft der Satz aus
dem Skript-Gerüst des Buchs: \enquote{Genau deshalb ist es ein
Experiment.} Für diese Einheit ist der Trick ohnehin nur Anlass; der Umbau
funktioniert auch, wenn der Text nur vorgelesen wird. Wer fragt, wie das
Seil wächst, bekommt die Gegenfrage, die zur Übung führt: \enquote{Was
genau haben Sie gesehen -- und was davon können Sie belegen?} Die Methode
wird nicht verraten (vgl.\ LE~13).
\end{vorbereitung}

\begin{ablauf}
\begin{description}
\item[0--10 · Staunmoment.] Vorführung mit dem blumigen Sprechtext. Danach
  liegt der Text auf der Folie; die Gruppe übersetzt ihn Satz für Satz in
  Fachsprache. Die Regeln, nach denen sie das tut, werden rechts an der
  Tafel gesammelt.
\item[10--40 · Theorie-Input.] IMRaD und Varianten; akademischer Stil und
  häufige Stilfehler; Argumentationsstruktur (These -- Beleg -- Schluss);
  Begriffspräzision am Beispiel NLP $\to$ evidenzbasiert; Zitieren;
  Plagiat, Paraphrase, Zitat.
\item[40--65 · Überarbeitungsübung.] Ein bewusst schlechter Absatz wird in
  zwei Runden überarbeitet: erst inhaltlich, dann sprachlich-formal
  (Übungsblatt, Aufgabe~2).
\item[65--85 · Eigene Einleitung.] Studierende schreiben die Einleitung
  ihrer Seminararbeit nach dem Vierschritt Thema -- Lücke -- Frage --
  Aufbau (Aufgabe~3).
\item[85--90 · Peer-Blick.] Partner-Feedback nach festem Raster
  (Aufgabe~4); Hinweis auf den Überarbeitungsprozess als Methode.
\end{description}
\end{ablauf}

%==============================================================================
\section{Staunmoment: Das wachsende Seil}
\label{sec:le11-patter}
%==============================================================================
\begin{staunbox}{Patter-Umbau}
Ein kurzes Seil, ein Sprechtext -- und am Ende ist das Seil deutlich länger.

\medskip
\emph{Hören Sie genau zu. Danach bekommen Sie den Text, und wir machen
daraus gemeinsam einen wissenschaftlichen Bericht.}
\end{staunbox}
\os{\newpage}

\nsl{Der folgende Sprechtext ist für diese Einheit geschrieben. Er sammelt
absichtlich die typischen Mittel eines Zauber-Patters: eine erfundene
Herkunftsgeschichte, vage Mengenangaben, eine unbelegte Ursache, einen
Appell an das Gefühl und eine Immunisierung gegen Zweifel.}
\begin{description}
\item[{\rot\bf Der Sprechtext (1/2)}]~\\[-\lblskip]
\begin{quote}\nsl{\small}
\enquote{Dieses kleine Seil hat mir ein alter Seemann in Lissabon
geschenkt. Er sagte, es habe ein Gedächtnis: Es erinnere sich an jede
Länge, die es je hatte. Sehen Sie, im Moment ist es kaum länger als meine
Hand. Aber wenn ich es behutsam durch die Finger gleiten lasse und ihm Zeit
gebe, dann erinnert es sich. Spüren Sie das?}
\end{quote}
\os{\newpage}

\item[{\rot\bf Der Sprechtext (2/2)}]~\\[-\lblskip]
\begin{quote}\nsl{\small}
\enquote{Es wächst, Zentimeter um Zentimeter, als würde es atmen. Kein
Mensch kann erklären, woher die Länge kommt; selbst die Wissenschaft steht
davor und staunt. Vielleicht ist es die Energie dieses Raumes, vielleicht
die Kraft Ihrer Aufmerksamkeit. Eines ist sicher: Je mehr Sie daran
glauben, desto länger wird es. Und jetzt schauen Sie -- so lang wie ein
Arm, und es hört nicht auf zu wachsen!}
\end{quote}
\os{\newpage}

\item[{\rot\bf Ihr Auftrag: übersetzen Sie}]~\\[-\lblskip]
\bi
  \ite Satz für Satz: Was \emph{behauptet} der Text?
  \ite Was davon haben Sie \emph{beobachtet}, was nur gehört?
  \ite Wie hieße der Satz in einem Beobachtungsbericht?
  \ite Notieren Sie zu jeder Änderung die \emph{Regel}, nach der Sie
    vorgegangen sind
\ei
\os{\newpage}
\end{description}

\nsl{Nach etwa fünf Minuten liegen an der Tafel die ersten Übersetzungen.
Eine mögliche Zielfassung, die die Lehrperson am Ende zeigt, lautet:}
\begin{description}
\item[{\rot\bf Eine mögliche Zielfassung}]~\\[-\lblskip]
\begin{quote}\nsl{\small}
Der Vortragende zeigte ein weißes Seil, dessen sichtbare Länge zu Beginn
auf etwa 15\,cm geschätzt wurde. Er ließ es etwa 30\,s lang durch beide
Hände laufen; danach war ein Seilstück von geschätzt 70\,cm sichtbar.
Gemessen wurde nicht; die Werte beruhen auf Schätzungen aus etwa 5\,m
Abstand. Ob die Materiallänge oder nur die sichtbare Länge zugenommen hat,
lässt sich aus diesen Beobachtungen nicht entscheiden. Die Herkunft des
Seils ist nicht belegt und für den Befund ohne Belang. Die vom Vortragenden
genannten Ursachen (\enquote{Energie des Raumes}, Aufmerksamkeit) wurden
nicht geprüft.
\end{quote}
\os{\newpage}
\end{description}

\begin{lehrhinweis}
Die Zahlen der Zielfassung sind Beispielwerte; die Gruppe soll die
\emph{eigenen} Schätzungen eintragen und merkt dabei, wie unterschiedlich
sie ausfallen. Das ist eine Brücke zu LE~9 (Intersubjektivität). Der
wichtigste Satz ist der über sichtbare Länge und Materiallänge: Er trennt
die Messgröße vom Konstrukt (LE~2) und deutet nur eine Richtung an, ohne
die Methode zu verraten. Das Buch nennt in der Praxis-Werkstatt eine
Fach-Variante (\enquote{Durch Zellteilung verlängert sich der
DNA-Strang}); sie zeigt den Registerwechsel, ist aber fachlich ungenau,
denn bei der Zellteilung wird die DNA verdoppelt, nicht verlängert. Wenn
Biologie-Studierende im Raum sind, lohnt der Hinweis: Fachsprache muss
auch fachlich stimmen.
\end{lehrhinweis}

%==============================================================================
\section{Vom Patter zur Fachsprache: die Transformationsregeln}
%==============================================================================
\nsl{Der Patter-Umbau zeigt, dass wissenschaftliche Sprache keine Frage des
Talents ist. Sie folgt Regeln, die man benennen und üben kann. Im Buch
geschieht dasselbe mit Begriffen aus dem Neuro-Linguistischen Programmieren
(NLP): Wo ein suggestiver Sammelbegriff wie \enquote{Anker} mehrere
getrennte Phänomene vermischt, setzt das Buch je einen Fachbegriff mit
Forschungsgrundlage ein; die Evidenzlage des NLP selbst ist dünn
\parencite{witkowski2010}. Präzision ist dabei eine Form epistemischer
Ehrlichkeit. Sie zwingt dazu, Behauptung und Beleg zu trennen und
Unsicherheit auszuweisen, statt sie rhetorisch zu überspielen.}
\begin{description}
\item[{\rot\bf Sieben Regeln des Umbaus}]~\\[-\lblskip]
\begin{center}\nsl{\small}
\begin{tabular}{p{0.36\textwidth}p{0.54\textwidth}}
\toprule
\textbf{Patter} & \textbf{Fachsprache} \\
\midrule
Erzählung (\enquote{ein Seemann sagte}) & Beobachtung: wer, was, wann, unter welchen Bedingungen \\
vage Menge (\enquote{Zentimeter um Zentimeter}) & Wert mit Einheit und Messverfahren \\
Behauptung ohne Beleg & Beleg angeben oder als ungeprüft kennzeichnen \\
Ursache (\enquote{Energie des Raumes}) & Hypothese mit Prüfbedingung \\
\bottomrule
\end{tabular}
\end{center}
\os{\newpage}

\item[{\rot\bf Sieben Regeln des Umbaus (Forts.)}]~\\[-\lblskip]
\begin{center}\nsl{\small}
\begin{tabular}{p{0.36\textwidth}p{0.54\textwidth}}
\toprule
\textbf{Patter} & \textbf{Fachsprache} \\
\midrule
Gefühlsappell (\enquote{Spüren Sie das?}) & Wirkung als Befund beschreiben (\enquote{Mehrere Zuschauende reagierten mit \ldots}) \\
Absolutheit (\enquote{kein Mensch}, \enquote{sicher}) & Geltungsbereich und Unsicherheit angeben \\
Immunisierung (\enquote{wenn Sie glauben}) & Bedingung nennen, unter der die Aussage scheitert \\
\bottomrule
\end{tabular}
\end{center}
\bi
  \ite Leitfrage für jeden Satz: \emph{Woher weiß ich das?}
\ei
\os{\newpage}

\item[{\rot\bf Begriffe ersetzen: das NLP-Beispiel}]~\\[-\lblskip]
\begin{center}\nsl{\small}
\begin{tabular}{p{0.3\textwidth}p{0.6\textwidth}}
\toprule
\textbf{statt (NLP)} & \textbf{besser (evidenzbasiert)} \\
\midrule
Anker setzen & Retrieval Cue \parencite{godden1975}; Implementation Intention \parencite{gollwitzer2006} \\
Reframing & Cognitive Reappraisal \parencite{gross1998} \\
Rapport aufbauen & psychologische Sicherheit \parencite{edmondson1999} \\
\bottomrule
\end{tabular}
\end{center}
\bi
  \ite ein Sammelbegriff verdeckt, \emph{welches} Phänomen gemeint ist
  \ite ein Fachbegriff verweist auf Befunde, die man nachlesen kann
\ei
\os{\newpage}
\end{description}

\begin{merkbox}[Präzision]
Präzise schreiben heißt nicht, kompliziert zu schreiben. Es heißt, bei
jedem Satz angeben zu können, woher man ihn weiß -- und kenntlich zu
machen, wo man es nicht weiß.
\end{merkbox}
\os{\newpage}

%==============================================================================
\section{Der Aufbau wissenschaftlicher Texte}
%==============================================================================
\nsl{Empirische Fachartikel folgen in den meisten Natur-, Ingenieur- und
Sozialwissenschaften dem Schema \emph{IMRaD}: Introduction, Methods,
Results and Discussion. Das Schema ist mehr als eine Konvention. Es trennt
die Teile eines Forschungsberichts nach ihrer Funktion: Was wollte ich
wissen? Wie habe ich es geprüft? Was kam heraus? Was bedeutet das? Wer
Ergebnis und Deutung in getrennten Kapiteln führt, kann die Deutung nicht
unbemerkt als Ergebnis ausgeben. Die Publikationsrichtlinien der American
Psychological Association legen dieses Schema für empirische Berichte
zugrunde \parencite{apa2020manual}.}
\begin{description}
\item[{\rot\bf IMRaD: vier Fragen, vier Teile}]~\\[-\lblskip]
\begin{center}
\resizebox{0.9\textwidth}{!}{%
\begin{tikzpicture}[font=\sffamily]
  \fill[rwuviolet!12] (0,0) -- (8,0) -- (5,-1.6) -- (3,-1.6) -- cycle;
  \draw[rwuviolet,thick] (0,0) -- (8,0) -- (5,-1.6) -- (3,-1.6) -- cycle;
  \node at (4,-0.7) {\textbf{Introduction}};
  \fill[rwucyan40!40] (3,-1.6) rectangle (5,-4.4);
  \draw[rwucyan,thick] (3,-1.6) rectangle (5,-4.4);
  \node at (4,-2.3) {\textbf{Methods}};
  \node at (4,-3.7) {\textbf{Results}};
  \draw[rwucyan,dashed] (3,-3) -- (5,-3);
  \fill[rwuviolet!12] (3,-4.4) -- (5,-4.4) -- (8,-6) -- (0,-6) -- cycle;
  \draw[rwuviolet,thick] (3,-4.4) -- (5,-4.4) -- (8,-6) -- (0,-6) -- cycle;
  \node at (4,-5.4) {\textbf{Discussion}};
  \node[right,font=\sffamily\small] at (8.4,-0.5) {Was wollte ich wissen? \emph{(allgemein $\to$ speziell)}};
  \node[right,font=\sffamily\small] at (8.4,-2.3) {Wie habe ich es geprüft?};
  \node[right,font=\sffamily\small] at (8.4,-3.7) {Was kam heraus? \emph{(ohne Deutung)}};
  \node[right,font=\sffamily\small] at (8.4,-5.5) {Was bedeutet das? \emph{(speziell $\to$ allgemein)}};
\end{tikzpicture}}
\end{center}
\os{\newpage}

\item[{\rot\bf Varianten}]~\\[-\lblskip]
\bi
  \ite \textbf{empirische Seminararbeit}: IMRaD, ergänzt um Abstract,
    Fazit, Literaturverzeichnis, Anhang (Material, Daten)
  \ite \textbf{Literatur- oder Theoriearbeit}: Einleitung -- thematisch
    gegliederter Hauptteil -- Diskussion -- Fazit
  \ite \textbf{technischer Bericht}: Anforderungen -- Entwurf --
    Umsetzung -- Test -- Bewertung
  \ite gemeinsam ist allen: Frage vorne, Antwort hinten, der Weg
    dazwischen nachvollziehbar
\ei
\os{\newpage}

\item[{\rot\bf Wohin gehört was?}]~\\[-\lblskip]
\bi
  \ite \enquote{Wir befragten 24 Studierende.} $\to$ Methode
  \ite \enquote{17 von 24 wählten Tür~B.} $\to$ Ergebnis
  \ite \enquote{Das deutet darauf hin, dass \ldots} $\to$ Diskussion
  \ite \enquote{Bisher ist unklar, ob \ldots} $\to$ Einleitung
  \ite häufigster Fehler: Deutung im Ergebnisteil
\ei
\nsl{Die Zahlen in diesen Beispielsätzen sind erfunden und dienen nur der
Zuordnung.}
\os{\newpage}
\end{description}

%==============================================================================
\section{Akademischer Stil}
%==============================================================================
\nsl{Akademischer Stil wird oft mit Unverständlichkeit verwechselt. Die
drei Anforderungen, auf die es ankommt, sind schlicht: \emph{Präzision}
(jeder Begriff hat eine Bedeutung, und zwar dieselbe im ganzen Text),
\emph{Sachlichkeit} (die Aussage stützt sich auf Belege, nicht auf
Begeisterung) und \emph{Transparenz} (Leserinnen und Leser können sehen,
woher eine Aussage stammt und wie sicher sie ist). Wer mit Fachvokabular
Eindruck machen will, erreicht oft das Gegenteil: In einer Reihe von
Experimenten beurteilten Versuchspersonen die Verfasser von Texten mit
unnötig komplizierten Wörtern als weniger intelligent als die Verfasser
einfacherer Fassungen desselben Inhalts \parencite{oppenheimer2006erudite}.}
\begin{description}
\item[{\rot\bf Drei Anforderungen}]~\\[-\lblskip]
\bi
  \ite \textbf{Präzision}: ein Begriff, eine Bedeutung -- im ganzen Text
  \ite \textbf{Sachlichkeit}: Belege statt Begeisterung
  \ite \textbf{Transparenz}: Herkunft und Sicherheit jeder Aussage sind
    erkennbar
  \ite einfach schreiben ist erlaubt: Imponiervokabular wirkt eher
    \emph{weniger} kompetent \parencite{oppenheimer2006erudite}
\ei
\os{\newpage}

\item[{\rot\bf Häufige Stilfehler}]~\\[-\lblskip]
\begin{center}\nsl{\small}
\begin{tabular}{p{0.3\textwidth}p{0.6\textwidth}}
\toprule
\textbf{Fehler} & \textbf{Beispiel} \\
\midrule
Nominalstil & \enquote{Die Durchführung der Befragung erfolgte \ldots} \\
Füllwörter & \enquote{eigentlich}, \enquote{quasi}, \enquote{ja}, \enquote{total} \\
Absolutheit & \enquote{Studien beweisen, dass \ldots} \\
Synonymwechsel & erst \enquote{Staunen}, dann \enquote{Verblüffung}, dann \enquote{Wonder} \\
unklarer Bezug & \enquote{Dies zeigt, dass \ldots} (was genau?) \\
\bottomrule
\end{tabular}
\end{center}
\os{\newpage}

\item[{\rot\bf Woher die Fehler kommen}]~\\[-\lblskip]
\bi
  \ite \textbf{inhaltliche Unsicherheit}: wer nicht weiß, was er
    behauptet, schreibt vage
  \ite \textbf{Imitation}: man kopiert den Klang von Fachtexten statt
    ihrer Logik
  \ite \textbf{kein Leserbild}: man schreibt für sich, nicht für jemanden,
    der prüfen will
  \ite Konsequenz: Stilprobleme sind oft Denkprobleme; man löst sie in
    Runde~1 der Überarbeitung
\ei
\os{\newpage}

\item[{\rot\bf Unsicherheit richtig dosieren}]~\\[-\lblskip]
\bi
  \ite zu sicher: \enquote{Staunen verbessert das Lernen.}
  \ite zu vorsichtig: \enquote{Es könnte möglicherweise eventuell sein,
    dass \ldots}
  \ite angemessen: \enquote{In der vorliegenden Stichprobe erinnerten
    Teilnehmende nach dem Staunmoment mehr Begriffe; ob der Effekt
    anhält, wurde nicht geprüft.}
  \ite Regel: so sicher formulieren, wie die Daten es tragen
\ei
\nsl{Auch der letzte Beispielsatz ist konstruiert und berichtet keinen
realen Befund.}
\os{\newpage}
\end{description}

%==============================================================================
\section{Argumentieren: These, Beleg, Schluss}
%==============================================================================
\nsl{Ein wissenschaftlicher Text ist eine Kette von Argumenten. Die
Grundform hat drei Teile: eine \emph{These} (was behauptet wird), einen
\emph{Beleg} (Daten, Quellen, Beobachtungen) und einen \emph{Schluss}, der
zeigt, warum der Beleg die These stützt. Stephen Toulmin hat dieses Schema
verfeinert. Zwischen Beleg und These steht bei ihm eine \emph{Schlussregel}
(warrant), die ihrerseits gestützt werden kann (backing); die These wird
durch einen \emph{Modaloperator} eingeschränkt (\enquote{vermutlich},
\enquote{in dieser Stichprobe}) und nennt die \emph{Ausnahmebedingungen},
unter denen sie nicht gilt \parencite{toulmin1958uses}. Die letzten beiden
Elemente sind genau das, was dem Patter fehlte.}
\begin{description}
\item[{\rot\bf Das Toulmin-Schema}]~\\[-\lblskip]
\begin{center}
\resizebox{0.95\textwidth}{!}{%
\begin{tikzpicture}[font=\sffamily,>=Stealth,
  b/.style={draw=rwuviolet,thick,fill=rwuviolet!8,rounded corners=2pt,
    minimum height=10mm,align=center,text width=34mm},
  c/.style={b,draw=rwucyan,fill=rwucyan40!40}]
  \node[b] (d) at (0,0) {\textbf{Beleg}\\Daten, Quelle};
  \node[c] (q) at (5.5,0) {\textbf{Modaloperator}\\\enquote{vermutlich}};
  \node[b] (t) at (11,0) {\textbf{These}\\Behauptung};
  \node[b] (w) at (2.75,-2.4) {\textbf{Schlussregel}\\warum Beleg $\to$ These};
  \node[b] (s) at (-2.6,-2.4) {\textbf{Stützung}\\Theorie, Studien};
  \node[c] (r) at (8.25,-2.4) {\textbf{Ausnahme}\\\enquote{außer wenn \ldots}};
  \draw[->,very thick] (d) -- (q);
  \draw[->,very thick] (q) -- (t);
  \draw[thick] (w.north) -- (2.75,0);
  \draw[->,thick] (s.east) -- (w.west);
  \draw[thick] (r.north) -- (8.25,0);
\end{tikzpicture}}
\end{center}
\os{\newpage}

\item[{\rot\bf Der Absatz als Argument}]~\\[-\lblskip]
\bi
  \ite \textbf{ein} Absatz = \textbf{ein} Gedanke
  \ite erster Satz: die These des Absatzes (Themensatz)
  \ite dann: Beleg mit Quelle, Erläuterung, ggf.\ Einschränkung
  \ite letzter Satz: was folgt daraus für die Frage der Arbeit?
  \ite Test: Lesen Sie nur die ersten Sätze aller Absätze. Ergibt sich
    der rote Faden?
\ei
\os{\newpage}
\end{description}

%==============================================================================
\section{Zitieren}
%==============================================================================
\nsl{Zitieren erfüllt drei Aufgaben: Es \emph{belegt} eine Aussage, es
\emph{würdigt} die Arbeit anderer, und es macht die Quelle
\emph{auffindbar}. Welche Form man wählt, ist zweitrangig, solange sie
vollständig und im ganzen Text einheitlich ist. Verbreitet sind drei
Familien: das Autor-Jahr-System (z.\,B.\ APA, 7.~Auflage
\parencite{apa2020manual}; auch dieses Skript zitiert so), das numerische
System (in Technik und Medizin verbreitet, z.\,B.\ IEEE oder Vancouver)
und das Fußnotensystem (Geisteswissenschaften, z.\,B.\ Chicago-Notes).
Maßgeblich ist, was Ihr Fach oder Ihre Prüferin vorgibt.
Literaturverwaltungsprogramme wie Zotero, Citavi oder JabRef nehmen die
Formatierung ab, die Verantwortung für die Richtigkeit der Angaben aber
nicht.}
\begin{description}
\item[{\rot\bf Wozu zitieren?}]~\\[-\lblskip]
\bi
  \ite \textbf{belegen}: Woher weiß ich das?
  \ite \textbf{würdigen}: Wessen Arbeit nutze ich?
  \ite \textbf{auffindbar machen}: Wo kann man nachprüfen?
  \ite Zitate sind die Fußspuren im Vertrauensspiel der Wissenschaft
    (\buch{Kap.~\enquote{Offenheit und Geheimnis}})
\ei
\os{\newpage}

\item[{\rot\bf Drei Zitierfamilien}]~\\[-\lblskip]
\begin{center}\nsl{\small}
\begin{tabular}{p{0.2\textwidth}p{0.36\textwidth}p{0.32\textwidth}}
\toprule
\textbf{System} & \textbf{im Text} & \textbf{verbreitet in} \\
\midrule
Autor-Jahr & (Toulmin, 1958, S.~xx) & Psychologie, Sozial-, Wirtschaftswiss. \\
numerisch & [12] & Technik, Informatik, Medizin \\
Fußnote & \textsuperscript{12} mit Vollangabe unten & Geschichte, Jura, Philosophie \\
\bottomrule
\end{tabular}
\end{center}
\bi
  \ite einheitlich bleiben -- und die Fachnorm erfragen
\ei
\os{\newpage}

\item[{\rot\bf Zitat und Paraphrase}]~\\[-\lblskip]
\bi
  \ite \textbf{direktes Zitat}: wörtlich, in Anführungszeichen, mit
    Seitenangabe; sparsam (Definitionen, prägnante Stellen)
  \ite \textbf{Paraphrase}: Gedanke in eigenen Worten \emph{und} eigener
    Satzstruktur -- braucht trotzdem einen Beleg
  \ite \textbf{Sekundärzitat} (\enquote{zit.\ nach}): nur, wenn das
    Original nicht erreichbar ist
  \ite Allgemeinwissen des Fachs braucht keinen Beleg -- im Zweifel
    belegen
\ei
\os{\newpage}
\end{description}

%==============================================================================
\section{Plagiat}
%==============================================================================
\nsl{Ein Plagiat liegt vor, wenn fremde Texte, Daten oder Ideen ohne
ausreichende Kennzeichnung als eigene ausgegeben werden. Die Absicht ist
für die Feststellung nicht entscheidend; auch nachlässiges Plagiieren ist
Plagiieren. Zusammen mit dem Erfinden und dem Verfälschen von Daten gehört
das Plagiat zu den drei Kerntatbeständen wissenschaftlichen Fehlverhaltens
(Fabrication, Falsification, Plagiarism, kurz FFP), die in LE~13
ausführlich behandelt werden. Für Prüfungsleistungen regelt die jeweilige
Prüfungsordnung die Folgen; in der Regel wird die Arbeit als
Täuschungsversuch gewertet. Ob und wie Texte aus KI-Werkzeugen verwendet und
gekennzeichnet werden dürfen, legen Hochschule und Prüfende fest; das
sollte vor Beginn der Arbeit geklärt sein.}
\begin{description}
\item[{\rot\bf Definition}]~\\
\defn{plagiat}{Ein \textbf{Plagiat} ist die Übernahme fremder Texte,
Daten, Abbildungen oder Ideen ohne ausreichende Kennzeichnung, sodass sie
als eigene Leistung erscheinen.}
\bi
  \ite entscheidend ist die \emph{Kennzeichnung}, nicht die Absicht
  \ite Plagiat ist eins von drei Kerndelikten (FFP $\to$ LE~13)
\ei
\os{\newpage}

\item[{\rot\bf Formen des Plagiats}]~\\[-\lblskip]
\bi
  \ite \textbf{Vollplagiat}: ganzer Text übernommen
  \ite \textbf{Copy \& Paste}: Passagen wörtlich ohne Anführungszeichen,
    auch wenn die Quelle irgendwo genannt ist
  \ite \textbf{verschleiernde Paraphrase}: Wörter ausgetauscht, Struktur
    und Gedanke übernommen, kein Beleg
  \ite \textbf{Übersetzungsplagiat}: fremdsprachige Quelle übersetzt, nicht
    belegt
  \ite \textbf{Selbstplagiat}: eigene, bereits bewertete Arbeit erneut
    eingereicht, ohne Hinweis
\ei
\os{\newpage}

\item[{\rot\bf Paraphrase oder Plagiat?}]~\\[-\lblskip]
\bi
  \ite Test 1: Könnte ich den Gedanken erklären, \emph{ohne} auf die
    Quelle zu schauen?
  \ite Test 2: Ist die Satzstruktur meine eigene?
  \ite Test 3: Steht der Beleg dort, wo der Gedanke steht -- nicht nur im
    Literaturverzeichnis?
  \ite dreimal \enquote{ja} $\to$ Paraphrase; sonst nachbessern
\ei
\os{\newpage}
\end{description}

%==============================================================================
\section{Einleitung und Fazit als Rahmen}
%==============================================================================
\nsl{Einleitung und Fazit sind die Teile, die am häufigsten gelesen werden,
oft als einzige. Für die Einleitung hat John Swales aus der Analyse vieler
Fachartikel ein Muster aus drei \enquote{Zügen} herausgearbeitet: das
Forschungsfeld abstecken, eine Lücke darin aufzeigen und die eigene Arbeit
in diese Lücke stellen \parencite{swales1990genre}. Der Vierschritt dieser
Einheit ist eine Lehrfassung dieses Musters. Das Fazit spiegelt die
Einleitung: Es beantwortet die Frage, die dort gestellt wurde, nennt die
Grenzen der Antwort und zeigt, welche Frage offen bleibt.}
\begin{description}
\item[{\rot\bf Die Einleitung im Vierschritt}]~\\[-\lblskip]
\begin{center}
\resizebox{0.95\textwidth}{!}{%
\begin{tikzpicture}[font=\sffamily,>=Stealth,
  st/.style={draw=rwuviolet,thick,fill=rwuviolet!10,rounded corners=3pt,
    minimum height=22mm,align=center,text width=36mm}]
  \node[st] (a) at (0,0) {\textbf{1 Thema}\\Worum geht es, warum ist es relevant?};
  \node[st] (b) at (4.6,0) {\textbf{2 Lücke}\\Was ist (belegt) noch unklar?};
  \node[st,draw=rwucyan,fill=rwucyan40!40] (c) at (9.2,0) {\textbf{3 Frage}\\Was genau untersuche ich?};
  \node[st] (d) at (13.8,0) {\textbf{4 Aufbau}\\Wie geht die Arbeit vor?};
  \draw[->,very thick] (a) -- (b);
  \draw[->,very thick] (b) -- (c);
  \draw[->,very thick] (c) -- (d);
\end{tikzpicture}}
\end{center}
\bi
  \ite Schritte 1--3: Feld abstecken, Lücke zeigen, Lücke besetzen
    \parencite{swales1990genre}
  \ite die Lücke braucht Belege -- sonst ist sie eine Behauptung
\ei
\os{\newpage}

\item[{\rot\bf Das Fazit spiegelt die Einleitung}]~\\[-\lblskip]
\bi
  \ite Frage aus der Einleitung wörtlich aufgreifen
  \ite Antwort in zwei, drei Sätzen -- ohne neue Daten
  \ite Grenzen der Antwort (Stichprobe, Methode, Geltungsbereich)
  \ite was offen bleibt: die nächste Frage
\ei
\os{\newpage}
\end{description}

%==============================================================================
\section{Überarbeiten als Methode}
%==============================================================================
\nsl{Kaum ein guter Fachtext entsteht in einem Durchgang. Erfahrene
Schreibende trennen das Schreiben vom Überarbeiten, und sie überarbeiten
in Runden mit unterschiedlichem Fokus. Die erste Runde fragt nach dem
Inhalt: Stimmt das Argument, tragen die Belege, ist die Reihenfolge
sinnvoll? Erst die zweite Runde gilt der Sprache und der Form. Wer beides
vermischt, poliert Sätze, die er danach streicht. Die Überarbeitung ist
damit das schriftliche Gegenstück zur Methode: Ein erster Entwurf ist eine
Hypothese über den Text, und die Überarbeitung prüft sie.}
\begin{description}
\item[{\rot\bf Zwei Runden}]~\\[-\lblskip]
\begin{center}\nsl{\small}
\begin{tabular}{p{0.44\textwidth}p{0.44\textwidth}}
\toprule
\textbf{Runde 1: Inhalt} & \textbf{Runde 2: Sprache und Form} \\
\midrule
Was ist die These? & Füllwörter und Absolutheiten streichen \\
Welche Belege fehlen? & Nominalstil auflösen \\
Was gehört nicht hierher? & Begriffe vereinheitlichen \\
Stimmt die Reihenfolge? & Zitate und Belege formal prüfen \\
\bottomrule
\end{tabular}
\end{center}
\bi
  \ite erst Runde 1 abschließen, dann Runde 2 beginnen
\ei
\os{\newpage}

\item[{\rot\bf Werkzeuge}]~\\[-\lblskip]
\bi
  \ite \textbf{Rückwärtsgliederung}: zu jedem Absatz einen Satz an den Rand
    schreiben -- ergibt die Satzfolge ein Argument?
  \ite \textbf{Abstand}: mindestens eine Nacht zwischen Schreiben und
    Überarbeiten
  \ite \textbf{laut lesen}: wo man stolpert, stolpert auch der Leser
  \ite \textbf{fremde Augen}: Peer-Feedback nach festem Raster
\ei
\os{\newpage}
\end{description}

\begin{merkbox}
Ein erster Entwurf ist eine Hypothese über den eigenen Text. Die
Überarbeitung prüft sie: erst den Inhalt, dann die Sprache.
\end{merkbox}
\os{\newpage}

%==============================================================================
\section{Übung in der Stunde}
%==============================================================================
\begin{description}
\item[{\rot\bf Überarbeitungsübung (Kleingruppen, 25 Min.)}]~\\[-\lblskip]
\bi
  \ite Absatz auf dem Übungsblatt (Aufgabe~2) lesen
  \ite \textbf{Runde 1} (rot, 10 Min.): Thesen, fehlende Belege,
    Überflüssiges markieren; neu ordnen
  \ite \textbf{Runde 2} (blau, 10 Min.): Sprache und Form
  \ite 5 Min.: Vergleich mit einer Nachbargruppe
\ei
\os{\newpage}

\item[{\rot\bf Ihre Einleitung (Einzelarbeit, 20 Min.)}]~\\[-\lblskip]
\bi
  \ite Vierschritt: Thema -- Lücke -- Frage -- Aufbau
  \ite Vorlage auf dem Übungsblatt (Aufgabe~3); heute ein Entwurf, als
    Hausaufgabe ca.\ 400 Wörter
  \ite danach 5 Min.\ Partner-Feedback nach Raster (Aufgabe~4)
\ei
\os{\newpage}
\end{description}

\begin{lehrhinweis}
Bei der Überarbeitungsübung neigen Gruppen dazu, sofort Sätze schöner zu
machen. Nach fünf Minuten von Runde~1 kurz durch die Gruppen gehen und
fragen: \enquote{Was ist die These dieses Absatzes?} Die meisten Gruppen
stellen fest, dass der Absatz mehrere Thesen hat und für keine einen Beleg.
Der Erwartungshorizont mit einer möglichen Endfassung steht am Ende des
Übungsblatts. Beim Peer-Feedback auf Ich-Botschaften achten
(\enquote{Mir war nach Schritt~2 nicht klar, \ldots}) und darauf, dass
Feedback beschreibt statt bewertet.
\end{lehrhinweis}

%==============================================================================
\section{Lernziele und Aufgaben}
%==============================================================================
\begin{description}
\item[{\rot\bf Nach dieser Einheit können Sie}]~\\[-\lblskip]
\bi
  \ite einen wissenschaftlichen Text nach IMRaD oder einer passenden
    Variante gliedern
  \ite Alltagssprache nach benennbaren Regeln in Fachsprache übertragen
  \ite ein Argument aus These, Beleg und Schluss aufbauen
  \ite Plagiat, Paraphrase und korrektes Zitat unterscheiden
  \ite einen Text in zwei Runden überarbeiten
\ei
\os{\newpage}

\item[{\rot\bf Bis zur nächsten Woche}]~\\[-\lblskip]
\bi
  \ite \textbf{Überarbeitung} des Absatzes abschließen (Aufgabe~2)
  \ite \textbf{Einleitung} Ihrer Seminararbeit, ca.\ 400 Wörter, nach dem
    Vierschritt (Aufgabe~3)
  \ite \textbf{Plagiats-Grenzfälle} einordnen (Aufgabe~5)
  \ite \textbf{Tagebuch des Staunens}: Welcher Ihrer Sätze hielt der Frage
    \enquote{Woher weiß ich das?} nicht stand?
  \ite für LE~12: 5-Minuten-Präsentation zum Semesterprojekt vorbereiten
\ei
\os{\newpage}
\end{description}

\handoutinput{handouts/LE11}
